% Contribution to GT_TOWER_CHECK; intended for the master's result environment.
% AI-produced by GPT-6 Astra (OpenAI, ChatGPT Pro mode). No priority claimed.
\subsection{The finite-subgroup theorem and the ten-point comparison}

Throughout this subsection a hat denotes the full profinite completion.
Let $\Gamma_{0,[n]}$ be the orientation-preserving mapping class group
of the sphere with $n\geq3$ punctures, allowing their permutation; its
pure subgroup is $\Gamma_{0,n}$.  A mixed group is a finite-index
subgroup of $\Gamma_{0,[n]}$.  All maps between discrete groups below
are homomorphisms; their completed maps are continuous homomorphisms.
Maps of fundamental groups induced without a chosen path between
basepoints are specified up to conjugation in their target.

\begin{result}{A.1.1}{Known}{The precise prime-power comparison}
For a mixed genus-zero mapping class group $M$ and a finite
$p$-subgroup $J\leq M$, the natural map $M\to\widehat M$ induces
a bijection on conjugacy classes of finite $p$-subgroups, and
\[
 N_{\widehat M}(J)=\overline{N_M(J)},\qquad
 C_{\widehat M}(J)=\overline{C_M(J)}.                    \tag{A.C1}
\]
This is Boggi--Zalesskii Theorem A \cite{CA_BZ2025}, equivalently the full
Theorem 4.3, applied to the completion map.  The centralizer assertion
is assertion (iii) of that theorem, not the last equality in the proof
of their Theorem 5.1(iii).
\end{result}
\begin{proof}
We check the reduction to the cited theorem.  The group
$\Gamma_{0,n}$ is the fundamental group of the ordered configuration
space of $n-3$ points in $\mathbb P^1\setminus\{0,1,\infty\}$.
Forgetting the last point gives successive extensions with finitely
generated free kernels.  The full profinite kernel topology and
exactness of the completed configuration sequences are
Boggi, \emph{Congruence topologies} \cite{CA_Boggi2020}, Corollary 4.7 and
Proposition 4.9, for the class of all finite groups.  Its
Corollary 5.4(ii) identifies full and congruence completion in genus
zero.  These statements pass to finite-index mixed groups by
restriction to the corresponding open subgroups.

Free groups are good.  The Lyndon--Hochschild--Serre comparison,
using the finite cohomology of finitely generated free groups with
finite coefficients, proves goodness inductively along these
extensions.  The configuration spaces have finite classifying
models, so their groups have finite cohomological dimension and type
$\mathrm{FP}_\infty$.  Finite extensions preserve goodness and
$\mathrm{FP}_\infty$ and give finite virtual cohomological dimension.
Thus $M$ is residually finite, good, and of finite virtual
$p$-cohomological type.

Theorem A is formulated for all discrete continuous
$\mathbb F_p[[\widehat M]]$-modules.  Such a module is a filtered
union of finite $\widehat M$-stable submodules: each vector has a
finite orbit and the span of finitely many finite orbits is finite.
Discrete group cohomology commutes with these unions because $M$ is
of type $\mathrm{FP}_\infty$.  Continuous cochains commute with them
because a continuous map from a compact space to a discrete module
has finite image.  The finite-module goodness comparison consequently
gives the module comparison required by Theorem A.  Its three
assertions give the stated conjugacy comparison and (A.C1).

There is also a direct way to read the final centralizer step in
\S4.3 of Boggi--Zalesskii.  Once their normalizer-density assertion
has been obtained, conjugation gives a continuous map
\[
 N_{\widehat M}(J)\longrightarrow\operatorname{Aut}(J).
\]
The target is finite.  Its kernel is therefore open, and the
intersection of the dense subgroup $N_M(J)$ with this kernel is
$C_M(J)$.  This proves the centralizer density without exchanging
closure with any intersection.
\end{proof}

\paragraph{The audited intersection step.}
For a finite subgroup $J=\langle j_1,\ldots,j_k\rangle$,
generation and the individual prime-power comparisons give
\[
 C_{\widehat M}(J)
   =\bigcap_i C_{\widehat M}(j_i)
   =\bigcap_i\overline{C_M(j_i)}.
                                                               \tag{A.C2}
\]
They do not supply the further equality with
$\overline{\bigcap_i C_M(j_i)}$.  That further equality is precisely
the simultaneous density still to prove.  It is asserted without
another argument in Boggi Proposition 2.4 and in the final part of
the proof of Boggi--Zalesskii Theorem 5.1(iii).  This audit identifies
a gap in those arguments for the general finite-subgroup assertion
\cite[Proposition 2.4]{CA_Boggi2026};
it supplies no counterexample to the assertions themselves.  For
the solvable genus-zero case we now give the comparison, including
the full subgroup topology, by a separate induction.

\begin{result}{A.1.2}{Proved}{Solvable finite subgroups in genus zero}
Let $M\leq\Gamma_{0,[n]}$ be of finite index.  The map
$M\to\widehat M$ induces a bijection between conjugacy classes of
finite solvable subgroups.  For every finite solvable $J\leq M$,
the subgroups $N_M(J)$ and $C_M(J)$ have the full profinite topology
induced from $M$, and
\[
 N_{\widehat M}(J)=\widehat{N_M(J)},\qquad
 C_{\widehat M}(J)=\widehat{C_M(J)}.                    \tag{A.C3}
\]
The identifications in (A.C3) are induced by the inclusions in $M$.
This proves the solvable comparison required in task A.1 and checks
LM4\_SUCCESSOR B.30/B.39 by an independent argument.
\end{result}
\begin{proof}
We first prove full induced topology for the centralizer and
normalizer of a finite $p$-group $A$.  This is a topology assertion
in addition to (A.C1).

Realize the finite group $A$ by orientation-preserving
homeomorphisms of the punctured sphere $S$.  Let $R$ be the set of
unmarked points with nontrivial stabilizer, put
$S^\circ=S\setminus R$, and let
$T^\circ=S^\circ/A$.  Then $T^\circ$ is a punctured sphere of
negative Euler characteristic, and the actual unramified cover gives
\[
 1\longrightarrow\Pi=\pi_1(S^\circ)
 \longrightarrow E=\pi_1(T^\circ)
 \longrightarrow A\longrightarrow1.                    \tag{A.C4}
\]
Both $\Pi$ and $E$ are nonabelian finitely generated free groups.
Write $\Gamma^\circ$ for the mixed mapping class group preserving
separately the old punctures and the newly marked set $R$.

The covering comparison has the exact sequences
\[
\begin{split}
 1&\longrightarrow A\longrightarrow N_{\Gamma^\circ}(A)
       \longrightarrow Q\longrightarrow1,\\
 1&\longrightarrow Z(A)\longrightarrow C_{\Gamma^\circ}(A)
       \longrightarrow Q_0\longrightarrow1.             \tag{A.C5}
\end{split}
\]
Here $Q$ is the subgroup of the downstairs mixed mapping class group
preserving the covering kernel $\Pi\triangleleft E$, with the
distinguished types of punctures, and $Q_0$ is the kernel of its
action on $A$ modulo inner automorphisms.  Each is of finite index:
there are finitely many homomorphisms from the finitely generated
group $E$ to the fixed finite group $A$.  The maps in (A.C5) are
descent through (A.C4), not identifications inferred from dimensions
of moduli spaces.  An upstairs normalizing mapping class descends;
a downstairs mapping class preserving the kernel lifts.  Changing
the lift multiplies by an element of $A$.  Requiring its action on
$A$ to be the identity gives the second exact sequence and kernel
$Z(A)$.  This is the ordinary covering mapping-class comparison
recorded explicitly in Boggi, \S2.1, equations (2)--(3)
of the published version \cite{CA_Boggi2026}.

For clarity, the free-group restriction used to pass to completions
is faithful in the required direction.  If an automorphism of
$\widehat E$ fixes its open normal subgroup $\widehat\Pi$
pointwise, then for every $e\in\widehat E$ its discrepancy
$e^{-1}\theta(e)$ centralizes $\widehat\Pi$.  That centralizer is
trivial: $\widehat\Pi$ contains nonzero integral powers of two
distinct free generators of $E$, whose procyclic centralizers have
trivial intersection.  Thus $\theta=1$.  If the restriction is
inner by an element of $\widehat\Pi$, we first compose with that
inner automorphism and apply the same argument.

Let $C^\circ=C_{\Gamma^\circ}(A)$ and intersect it with the pure
mapping class group on all punctures of $S^\circ$.  The resulting
$C_1$ is induced-open in $C^\circ$ and meets $Z(A)$ trivially:
a finite orientation-preserving sphere action fixing three labelled
points is the identity.  Under (A.C5), $C_1$ identifies with a
finite-index subgroup of $Q_0$.  Genus-zero congruence faithfulness
on $T^\circ$ makes the completed action of $\widehat C_1$ on
$\widehat E$ faithful.  The preceding restriction argument makes
its action on $\widehat\Pi$ faithful.  The action factors through
$\widehat{\Gamma^\circ}$, so
$\widehat C_1\to\widehat{\Gamma^\circ}$ is injective.  Since
$C_1$ is induced-open, the kernel of
$\widehat C^\circ\to\widehat{\Gamma^\circ}$ lies in
$\widehat C_1$ and is zero.  This proves full induced topology for
$C^\circ$.

Filling $R$ has the exact completed configuration sequence
\[
 1\longrightarrow\widehat K\longrightarrow
 \widehat{\Gamma^\circ}\longrightarrow\widehat{\Gamma(S)}
 \longrightarrow1,                                      \tag{A.C6}
\]
where $K=\pi_1(\operatorname{UConf}_{|R|}(S))$ for the unordered
new marks.  Exactness and full kernel topology follow by the
configuration result already cited, followed by the finite
permutation extension.  Discretely, filling gives isomorphisms
of the two normalizers and of the two centralizers of $A$.
Indeed the downstairs groups and both kernels in (A.C5) are
unchanged: the additional punctures are exactly the branch points
already specified by the covering.  Equivalently these are the two
commutative exact-row diagrams of Boggi Proposition 2.1.  In
particular $A\cap K=1$ and $C_K(A)=1$.

The subgroup $K\cdot A=K\rtimes A$ is a finite extension of a
configuration group and hence is residually finite, good and of
finite virtual cohomological type by the free-kernel argument used
above.  Its completion is the actual subgroup
$\widehat K\rtimes A$ in (A.C6), because $K$ has full induced
topology there and the finite quotient $A$ is detected in the base.
Direct application of Theorem A(iii) to this single $p$-group gives
\[
 C_{\widehat K\rtimes A}(A)
   =\overline{C_{K\rtimes A}(A)}=Z(A),\qquad
 C_{\widehat K}(A)=1.                                    \tag{A.C7}
\]
Thus the embedded $\widehat C^\circ$ has trivial intersection with
the kernel in (A.C6), and its image after filling is the full
completion of $C_{\Gamma(S)}(A)$.  This is the prime-power repair
of the one centralizer-density input in Boggi \S3.2; no general
finite-subgroup intersection assertion is used.

The centralizer is finite index in the normalizer.  It is also
induced-open: choose a finite-index normal subgroup
$U\triangleleft\Gamma(S)$ with $U\cap A=1$.  For
$g\in N_{\Gamma(S)}(A)\cap U$, every commutator $[g,a]$
belongs to $U\cap A$, hence vanishes.  Consequently a finite
extension transfers the full induced topology to the normalizer.
Intersecting with an arbitrary finite-index mixed group $M$
preserves all these topology assertions.  Together with (A.C1),
we have proved
\[
 N_{\widehat M}(A)=\widehat{N_M(A)}
                                                        \tag{A.C8}
\]
for every finite $p$-subgroup $A$.

We now induct on the order of a solvable finite subgroup.  The
three assertions to be proved together are realization of all
finite solvable subgroups of $\widehat M$, injectivity on their
discrete conjugacy classes, and (A.C3) with full topology.  They
are immediate for the trivial group.  In a nontrivial finite
solvable group $J$, let $D$ be the last nontrivial derived group,
choose a prime $p\mid|D|$, and set $A=D[p]$.  This is a nontrivial
characteristic elementary abelian $p$-subgroup.  The same choice is
preserved by every isomorphism of $J$.

First suppose $J\leq M$.  Put $N=N_M(A)$ and $Q=N/A$.
Descent through the actual $A$-cover, with all unmarked branch
points punctured, identifies $Q$ with a finite-index mixed
genus-zero mapping class group by the first sequence in (A.C5).
Its full completion is $\widehat Q=\widehat N/A$; the finite
kernel is closed and survives because $N$ is residually finite.
Characteristicity of $A$ and (A.C8) give exact preimage equalities
\[
\begin{split}
 N_M(J)&=\operatorname{pr}^{-1}N_Q(J/A),\\
 N_{\widehat M}(J)&=\widehat{\operatorname{pr}}^{-1}
                         N_{\widehat Q}(J/A),
\end{split}                                                  \tag{A.C9}
\]
inside $N$ and $\widehat N$, respectively.  Indeed $J$ is the
full inverse image of $J/A$ in either group.  Induction applies
to $J/A$ in $Q$.  Taking the inverse image of its dense discrete
normalizer proves density in the second line of (A.C9).

The topology transfer in this step can be made explicit.  Choose
a finite-index subgroup $V\leq N$ avoiding $A$.  Its image
$\overline V=VA/A$ is finite index in $Q$, and $V\to\overline V$
is an isomorphism.  If $L$ is either normalizer in (A.C9),
the inclusion $L\cap V\hookrightarrow V$ identifies with
$\operatorname{pr}(L)\cap\overline V\hookrightarrow\overline V$.
Induction gives the full induced topology for that inclusion.
Every finite-index subgroup of $V$ is also finite index in $N$,
so its core gives the required finite-index normal subgroup of
$N$.  The finite-index extension from $L\cap V$ to $L$ proves
the topology assertion for $N_M(J)$, and (A.C8) transfers it to
$M$.  Finally the open kernel of
$N_{\widehat M}(J)\to\operatorname{Aut}(J)$ gives centralizer
density.  The centralizer is induced-open by the same finite
quotient avoiding $J$, so it too has full induced topology.

For realization, take a finite solvable $J\leq\widehat M$.
Theorem A conjugates its chosen $A$ to a discrete $p$-subgroup.
Now $J$ is contained in the normalizer (A.C8).  Induction
conjugates $J/A$ inside $\widehat Q$ to a finite subgroup
$\overline J_0\leq Q$.  Lift that conjugator to $\widehat N$.
The full inverse image $J_0$ of $\overline J_0$ in $N$ is
finite, of order $|A|\,|\overline J_0|$, and its inverse image
in $\widehat N$ is the same finite subgroup.  The conjugated
$J$ is exactly $J_0$.

For injectivity, let discrete solvable subgroups $J_1,J_2$ be
conjugate in $\widehat M$.  The conjugacy identifies their
chosen characteristic $p$-subgroups.  The prime-power injectivity
of Theorem A first conjugates these discretely to the same $A$.
The remaining profinite conjugator lies in $\widehat N$.
Inductive injectivity in $Q$ supplies a discrete conjugator of
$J_1/A$ and $J_2/A$; any lift to $N$ conjugates their full
preimages $J_1,J_2$.  This completes the induction and the proof.
\end{proof}

\begin{result}{A.1}{Proved}{Audit of EXIT\_TOWER E.44}
The ten-point argument of EXIT\_TOWER E.44 proves
$\mathcal J(A)=A$ for every $A\in\widehat{\mathrm{GT}}$.
Its finite-subgroup input is Theorem A(iii) applied directly to
$H=C_2^2$ in a finite-index subgroup of $\Gamma_{0,[10]}$.
It does not use the closure--intersection equality (A.C2).
The centralizer has its full induced topology by an explicit
retraction.  Thus Lead 1(3) is confirmed.  The literal explanation
in Lead 2 is corrected: the relevant prime-power theorem already
supplies the simultaneous centralizer, while its separate solvable
claim is proved by A.1.2.
\end{result}
\begin{proof}
We give the geometric and outer-naturality comparisons; they use
only the ordinary genus-zero GT tower action.  That action on all
full completed sphere mapping class groups, compatible with
forgetting and relabelling, is the usual consequence of Drinfeld's
relations; see the standard action (1.1) and its passage to mixed
sphere groups in Nakamura--Schneps, \S1 \cite{CA_NS2000}.
The pentagon enters at this stated input.  No action on
a higher-genus completion is used to prove the identity below.

On $\mathbb P^1$ use the commuting involutions
$a(w)=-w$, $b(w)=1/w$, and put $c=ab$.  Mark
\[
\begin{gathered}
 H=\langle a,b\rangle=C_2^2,\quad
 A_0=(0,\infty),\quad B_0=(1,-1),\quad C_0=(i,-i),\\
 D_0=(z,-z,1/z,-1/z),\qquad
 z\in\mathbb P^1\setminus\{0,\infty,\pm1,\pm i\}.
\end{gathered}                                                \tag{A.C10}
\]
The subscripts on these ordered sets distinguish them from the GT
parameter $A$.  Let $G$ be the inverse image of
$C_{\mathfrak S_{10}}(H)$ in $\Gamma_{0,[10]}$, and set
$N=C_G(H)$.  The indicated permutation centralizer has order
$2^3\cdot4=32$: it flips each of the three fixed pairs and
translates the regular $H$-set $D_0$.

Removing all ten marks and all four quotient marks gives the
actual unramified cover
\[
 1\longrightarrow F_9\longrightarrow F_3
       \longrightarrow H\longrightarrow1.                 \tag{A.C11}
\]
The three branch peripheral elements map to $a,b,c$ and the
fourth to $1$.  Conjugation identifies the middle free group
with the inverse image of $H$ in $\operatorname{Aut}(F_9)$.
Indeed its centralizer of $F_9$ is trivial, and the kernel over
$H$ is exactly the inner automorphism group of $F_9$.
An element of $N$ therefore descends to an outer automorphism
of the middle group acting trivially on $H$.

At a branch puncture the square of a positive quotient
peripheral is an upstairs peripheral.  The descended image of
that square is conjugate to the same positive square.  Unique
roots in the discrete free group imply the corresponding
assertion for the peripheral itself.  The fourth puncture
requires no square.  Punctured-sphere Dehn--Nielsen--Baer
therefore identifies the descended class with a pure
four-point mapping class.  Conversely every pure downstairs
class fixes the homomorphism to the abelian group $H$, hence
lifts and centralizes $H$.  The deck kernel is $H$.  This proves
$N/H\cong\Gamma_{0,4}$ by the actual covering map.

Forgetting $A_0,B_0,C_0$ gives the stronger isomorphism
\[
 f_4:N\xrightarrow{\sim}\Gamma_{0,4}\times H_D,
 \qquad N/H\xrightarrow{\sim}\Gamma_{0,4},                 \tag{A.C12}
\]
where $H_D$ is the canonical Klein four group of double
transpositions of $D_0$.  To check the asserted map, put
\[
 u=\frac{(z^2+1)^2}{4z^2},\qquad
 \lambda_D=\left(\frac{z^2-1}{z^2+1}\right)^2=1-\frac1u.
                                                               \tag{A.C13}
\]
The first is the quotient coordinate; the second is the
cross-ratio of the ordered quartet $D_0$.  Their relation is
an isomorphism of the thrice-punctured sphere, and the deck
map is $H\cong H_D$.  Thus the maps on both ends of the exact
centralizer sequences are isomorphisms, proving (A.C12).
One can also invert the family: take the canonical Klein four
group of a quartet and add its three fixed pairs.  They are
distinct and disjoint from the quartet, and this construction
is inverse on equivariant isomorphisms.

Every element of $G$ has its $D_0$ permutation in $H_D$, so
forgetting is defined on all of $G$ with the same target as
(A.C12).  Its composition with $f_4^{-1}$ is a retraction
$G\to N$.  It follows that $\widehat N\to\widehat G$ is
injective and has image $\overline N$.  Theorem A(iii), for
the one group $H$, now gives
\[
                  C_{\widehat G}(H)=\widehat N.             \tag{A.C14}
\]

There are exactly two pure conjugacy classes for the
elementwise label permutations (A.C10).  Finite realization on
the punctured sphere \cite[Theorem 5 and \S IV, pp.~255--256]{CA_Kerckhoff1983},
followed by uniformization, permits the
normalization $A_{0,1}=0$, $A_{0,2}=\infty$, $B_{0,1}=1$.
The two involutions then have the displayed formulas, so
$B_{0,2}=-1$ and the two possible $C_0$ orders are
$(i,-i)$ and $(-i,i)$.  The allowed domain of $z$ is connected;
each of these orders therefore gives one pure conjugacy class.
They are distinct also by the following finite monodromy check.
For the canonical Klein four action on a normalized quartet,
the fixed-pair discriminants are
$4\lambda_D$, $4(1-\lambda_D)$ and
$4\lambda_D(\lambda_D-1)$.  Loops around $0,1$ flip
$(A_0,C_0)$ and $(B_0,C_0)$, respectively.  Together with
the four deck translations, this gives permutation image of
$N$ of order $16$, of index two in the order-$32$ permutation
centralizer.  Within $G$ every normalizer of $H$ centralizes
it, since its conjugation acts trivially on the faithful
permutations of $H$.  The two pure double cosets are exactly
the two classes just described.

Let the ordinary genus-zero tower action associated to
$A\in\widehat{\mathrm{GT}}$ act on $\widehat G$.  It induces
the identity on the finite label-permutation quotient.
Theorem A(i) conjugates its image of $H$ to a discrete group.
Choose a discrete element of $G$ with the same finite
permutation as that conjugator and subtract it; the remaining
conjugator is pure, while the discrete subgroup changes by
that discrete conjugation.  The classification then requires
only a further pure conjugation and possibly
$P=(C_{0,1}\ C_{0,2})$.  After this correction the action
fixes $H$ pointwise, since its faithful label permutations
are unchanged.  Represent $P$ by a half-twist in a disk
containing just its two labels.  It is killed by forgetting
$C_0$, hence by every forgetting map used below.  The pure
corrections introduce only the target inner conjugacies
associated with those maps.  Both pure classes have thus been
accounted for.

The target is the mixed four-point orbifold group
\[
 R_0=\Gamma_{0,4}^{D_8}/V
       \cong\langle v,d\mid(vd)^2=1\rangle\cong C_2*\mathbb Z,
 \qquad R=\widehat R_0.                                  \tag{A.C15}
\]
Here $V$ is the canonical normal Klein four group, $D_8$
preserves an unordered partition of the quartet into two
blocks, and $C_4=\mathfrak S_2\times\mathfrak S_2$
preserves the two blocks separately.  The actual colored
forgetting targets suffice because
\[
 C_4V=D_8,\qquad C_4\cap V=C_2,
 \qquad \Gamma_{0,4}^{C_4}/C_2\xrightarrow{\sim}R_0.      \tag{A.C16}
\]
These are homomorphisms of the actual mixed groups and finite
quotients, so they persist under full completion.

The quotient by $V$ is preserved by the tower action also
after completion.  Indeed any finite normal subgroup of the
completed full four-point group intersects its torsion-free
pure $\widehat F_2$ trivially and hence centralizes that pure
group.  The quotient $\mathfrak S_4/V=\mathfrak S_3$ acts
faithfully on its three peripheral homology vectors.
Consequently its centralizer is contained in
$\widehat F_2V$ and then, by centerlessness, equals $V$.
Thus $V$ is the maximal finite normal subgroup and is
characteristic.

Let $N_b$ stabilize $O_b=\{z,1/z\}\subset D_0$.
There are two $b$-orbits in $D_0$, interchanged by $H$;
hence
\[
 N=HN_b,\quad H\cap N_b=\langle b\rangle,
 \quad N_b/\langle b\rangle\xrightarrow{\sim}N/H.         \tag{A.C17}
\]
Invert this last isomorphism, retain $A_0\sqcup O_b$ by
forgetting, and use (A.C16).  The resulting homomorphism
$q_b:N/H\to R_0$ is well-defined because $b$ is the double
flip killed in that target.  Similarly, the stabilizer of
$O_a=\{z,-z\}$, forgetting to $B_0\sqcup O_a$, and killing
the double flip $a$, gives $q_a:N/H\to R_0$.
Both stabilizers have index two.  Thus the same construction
on completions gives continuous $q_a,q_b:\widehat{N/H}\to R$.

Use the same ordered cross-ratio convention
\[
 \operatorname{cr}(p_1,p_2;p_3,p_4)
 =\frac{(p_3-p_1)(p_4-p_2)}{(p_3-p_2)(p_4-p_1)}
\]
for both retained quartets.  It gives
\[
\begin{aligned}
 \ell_b&=\operatorname{cr}(0,\infty;z,1/z)=z^2,
 &U_b&=\frac{(\ell_b+1)^2}{4\ell_b}=u,\\
 \ell_a&=\operatorname{cr}(1,-1;z,-z)
                  =\left(\frac{z-1}{z+1}\right)^2,
 &U_a&=\frac{(\ell_a+1)^2}{4\ell_a}=\frac{u}{u-1}.
\end{aligned}                                                \tag{A.C18}
\]
The target orbifold has a cone of order two at $U=0$
and punctures at $1,\infty$.  Therefore
\[
 q_a=S_Rq_b,\qquad S(U)=\frac U{U-1},\qquad
 S_R(v)=d,\quad S_R(d)=v                              \tag{A.C19}
\]
as outer homomorphisms, with one ordinary target basepoint
conjugacy.  We choose that comparison once so that the first
identity is literal.  In this marking $q_b$ is the surjection
\[
 \widehat F(v,d)\longrightarrow
       \langle v,d\mid(vd)^2=1\rangle^{\wedge}.             \tag{A.C20}
\]
It fills the square of the peripheral at the missing cone
point.  Its kernel is the closed normal subgroup generated
by $(vd)^2$.  Surjectivity follows from the displayed discrete
presentation and right exactness of full completion.

The colored target's pure subgroup is
\[
 K_{\rm std}=\ker(R\to C_2),\quad v\mapsto0,
 \ d\mapsto1,\qquad K_{\rm std}=\widehat F(v,d^2).        \tag{A.C21}
\]
This cusp assignment can be read directly from
$U=(\ell+1)^2/(4\ell)$: $U=1$ has the ramified lift
$\ell=1$, whereas $U=\infty$ has the two degree-one
lifts $\ell=0,\infty$.  Thus its standard ordered
peripheral generators are $v,d^2$.  The same convention
is used for both quartets in (A.C18); there is no swap of
a cone and a puncture, or an unrecorded power of a cusp.
Changing the pure quartet coordinate by a permutation of
$0,1,\infty$ conjugates by its geometric
$\mathfrak S_3$ action.  The GT outer action commutes with
that action, by its extension to the mixed four-point
group (equivalently the two-cycle and three-cycle
relations).  Thus (A.C13) also identifies the source
outer parameter with the same named $A$.

Let $\iota(A)\in\operatorname{Out}(R)$ denote the extension
of the standard action on (A.C21).  Its representative is
\[
 v\longmapsto v^\lambda,\qquad
 d\longmapsto f(v,d^2)^{-1}d^\lambda f(v,d^2)
 \quad\text{for }A=(\lambda,f).                          \tag{A.C22}
\]
Existence is the ordinary mixed four-point tower action.
The extension is unique in the required outer sense:
align two restrictions by an inner element of $K_{\rm std}$;
their quotient fixes $K_{\rm std}$ pointwise.  Its
discrepancy at any $r\in R$ centralizes that open normal
subgroup, whose centralizer is trivial.  Hence the
quotient is the identity.  To check injectivity including
the possible deck inner automorphism, suppose
$\iota(A)=\operatorname{Inn}(h)$ in $R$.
The normalized GT action on $K_{\rm std}^{\rm ab}$ is
the scalar matrix $\lambda I$.  Conjugation by the
nontrivial coset $dK_{\rm std}$ has matrix
$\left(\begin{smallmatrix}-1&0\\-1&1\end{smallmatrix}\right)$
in the basis $(v,d^2)$, by $dvd^{-1}=v^{-1}d^{-2}$.
It is not scalar.  Thus $h\in K_{\rm std}$, and
restriction is inner by a pure element.  The normalized
GT action on $\widehat F_2$ is faithful as an outer
action, so $A=1$.  This proves injectivity with the
deck ambiguity removed.

After the finite-subgroup correction above, the action
preserves $H$, $\widehat N$, both choice subgroups, and
the displayed forgetting maps and finite quotient kernels.
The source action $\alpha$ on $\widehat N/H$ has outer
parameter $A$ by (A.C12), and each colored target action
$\beta$ has outer class $\iota(A)$ by (A.C21).
There are consequently elements $g_a,g_b\in R$ with
\[
 q_b\alpha=\operatorname{Inn}(g_b)\beta q_b,
 \qquad q_a\alpha=\operatorname{Inn}(g_a)\beta q_a.
                                                               \tag{A.C23}
\]
These inner elements are allowed to differ.  Substituting
(A.C19) into (A.C23) and cancelling the common
surjection (A.C20) gives
\[
 S_R\beta
  =\operatorname{Inn}\bigl(S_R(g_b)^{-1}g_a\bigr)\beta S_R.
                                                               \tag{A.C24}
\]
Thus $[\beta]$ commutes with $[S_R]$ in
$\operatorname{Out}(R)$.  EXIT\_TOWER E.28 defines the
other reading by restriction to
$K_{\rm geo}=\widehat F(d,v^2)$ and proves the explicit
comparison
\[
 \iota(\mathcal J(A))=[S_R]\iota(A)[S_R].                 \tag{A.C25}
\]
Equations (A.C24)--(A.C25) and injectivity of $\iota$
prove $\mathcal J(A)=A$.

Finally, the preliminary branch-order recognition in E.44
uses the actual universal-curve preimage
$C_2*C_2*C_2*C_2$, with the even kernel the twice-punctured
torus group.  Discrete reduced words give four conjugacy
classes of order-two lifts; direct Theorem A at $p=2$
gives the same after full completion.  Each is one
even-kernel class because its odd generator centralizes
it.  Their permutation is the four-strand permutation;
ordering them gives the index-$24$ cover and the named
five-point kernel.  This recognizes that cover, while
(A.C10)--(A.C25) supplies its parameter comparison.
The recognized cover therefore agrees with the named
E.37 construction, whenever interpreting it in that
already established ordinary universal-curve diagram.
The proof of (A.C24), and hence of the genus-zero
identity itself, uses no higher-genus tower extension.
\end{proof}

\paragraph{Executed finite checks and their exact scope.}
The independent program \texttt{check\_e44\_geometry.py}
checks the order-$4$ group, its order-$32$ permutation
centralizer, the order-$16$ geometric image and two
double cosets, both choice subgroups, their intersections
with $H$, the killed correction $P$, (A.C16), the
rational identities (A.C13) and (A.C18), and the three
fixed-pair discriminants.  The group-theoretic topology
and outer-naturality arguments (A.C14), (A.C23) and
(A.C24) are proved above; they are not inferred from
these finite computations.
